% =============================================================================
%  CONCLUSION GÉNÉRALE ET PERSPECTIVES
% =============================================================================
\cleardoublepage
\entete{Conclusion générale}
\chapter*{CONCLUSION GÉNÉRALE}
\addcontentsline{toc}{chapter}{CONCLUSION GÉNÉRALE}

Au terme de ce travail, mené du 20 juin au 20 août 2026 au sein du \emph{Think
Tank} Fou'ndjem de la Fondation Jean Félicien Gacha, le point de départ mérite
d'être rappelé : la découverte de la Fondation supposait de s'y rendre, ce qui
écartait les publics éloignés, la diaspora et les personnes à mobilité réduite,
et là où une présence numérique existait, elle se réduisait à une vitrine figée.
La question posée était donc celle d'une plateforme unifiée capable d'offrir une
visite en trois dimensions, de retrouver les œuvres apparentées dans d'autres
institutions et d'assurer une assistance fiable fondée sur les seules données de
la Fondation.

Pour y répondre, nous avons conduit une recherche interactive de niveau
intégrateur, selon une approche mixte associant l'analyse documentaire,
l'observation directe, six entretiens semi-directifs et une enquête auprès de
quarante-deux visiteurs. Le travail a abouti à une plateforme opérationnelle,
MUSÉA, articulant vingt fonctionnalités autour de quatre ensembles : un
progiciel de gestion couplé à un site public éditable ; une visite immersive à
360~degrés complétée par la 3D et la réalité augmentée à taille réelle ; un
assistant conversationnel ancré sur les documents validés ; et un dispositif de
rapprochement avec les collections ouvertes du monde. Les quatre hypothèses
spécifiques ont été vérifiées et l'hypothèse générale corroborée.

Au-delà de l'outil livré, trois enseignements constituent la contribution propre
de ce travail. \textbf{Le décalage de vocabulaire constitue d'abord une seconde
dispersion des œuvres} : une recherche en français dans les catalogues
internationaux ramène du bruit là où la même recherche, reformulée dans leur
vocabulaire de catalogage, identifie des correspondances précises, le meilleur
score passant de 57 à 85 sur 100 ; à la dispersion physique s'ajoute donc une
dispersion documentaire, réversible celle-là. \textbf{L'isolement des données ne
peut ensuite pas être une propriété du code} : les deux fuites détectées étaient
d'origine applicative, et seule une garantie portée par la base de données
résiste à l'erreur de programmation. \textbf{La qualité d'une visite restituée
repose enfin sur quatre continuités}, spatiale, d'échelle, documentaire et
narrative, dont la rupture d'une seule fait échouer l'ensemble.

Sur le plan personnel, ce stage aura permis de conduire un projet complet, du
besoin exprimé jusqu'à la mise en production. Mené auprès d'une institution
unique, ce travail ne prétend pas à la généralisation : sa transposition
demandera vérification, objet des perspectives ci-après.


% =============================================================================
\section*{Perspectives}
\addcontentsline{toc}{section}{Perspectives}

Le dispositif construit ouvre des prolongements que nous regroupons en quatre
axes, du plus immédiat au plus prospectif. Certains n'ont pu être menés à terme
durant le stage, faute de temps, de moyens ou d'accès ; nous le signalons.

\textbf{Consolider ce qui existe.} La perspective la plus importante est de
comparer les œuvres sur l'image et non sur le texte : les rapprochements
reposent aujourd'hui sur la description écrite, faute d'avoir pu installer un
modèle de vision sur le poste de travail. Un plongement visuel rapprocherait
deux objets semblables mais décrits selon des conventions différentes, situation
fréquente entre un catalogue africain et un catalogue européen ; il supprimerait
la principale limite de la mémoire réunifiée. Viennent ensuite la production des
cartes de profondeur sur des œuvres réelles, que l'exécution en conteneur
rendrait possible, l'obtention des clés d'accès aux quatre collections écartées,
qui ne dépend que d'une démarche administrative, et l'épreuve de la chaîne de
rapprochement de bout en bout, chaque maillon n'ayant été vérifié qu'isolément.

\textbf{Enrichir l'expérience du visiteur et l'ancrer dans la culture locale.}
L'écran de fin de parcours existe ; il reste à y animer le globe de la
dispersion, pour que la visite s'achève sur la mesure de ce dont l'institution a
été séparée. Le mode hors connexion, en mettant en cache les panoramas des
salles visitées, permettrait de poursuivre une visite sur une connexion
défaillante, situation courante pour une partie du public visé. Surtout, l'ajout
des langues de l'Ouest camerounais, ghomálá', fe'efe'e et medumba, à une
application aujourd'hui bilingue transformerait la portée du projet : il ne
s'agirait plus seulement de faire connaître un patrimoine à l'extérieur, mais de
le restituer aux communautés qui en sont issues, dans leur langue, ce que la
structure du système permet sans modification. La collecte des récits oraux
auprès des détenteurs de la mémoire, rattachée aux fiches d'œuvres et aux
personnages de la généalogie, enrichirait enfin la base documentaire de
l'assistant.

\textbf{Industrialiser, pérenniser et essaimer.} Le passage aux conteneurs
orchestrés, décrit au chapitre~4, conditionne l'automatisation de l'ingestion
documentaire, puis la migration vers Kubernetes au-delà d'une dizaine
d'institutions. L'accueil d'institutions partenaires et l'ouverture des données
de la Fondation, selon les mêmes standards que les catalogues qu'elle interroge,
rendraient le mouvement réciproque. Enfin, la mutualisation du catalogue mondial
fait décroître le coût par institution à mesure que le réseau s'étoffe : de
petites structures accèdent ainsi à un outil qu'elles ne pourraient financer
isolément.
